% Part: sets-functions-relations
% Chapter: infinite
% Section: dedekinds-proof
%
\documentclass[../../../include/open-logic-section]{subfiles}

\begin{document}

\olfileid{sfr}{infinite}{dedekindsproof}

\olsection[Dedekind's ``Proof'']{Dedekind 对无限集存在性的``证明''}

在本章中，我们借助 Dedekind 代数对自然数给出了一种集合论处理。在\olref[arith][ref]{sec}，我们已反思过分析的算术化（及其他问题）的哲学意义。现在，我们应该反思我们在本章所做工作的意义。

贯穿\olref[sfr][arith][]{chap}，我们将自然数视为给定的，并明确地用它们来构造整数、有理数和实数。在本章，我们并未给出自然数的一个显式构造。我们仅仅证明了，\emph{给定任一 Dedekind 无限集}，我们可以定义一个集合，它的行为将完全符合我们对~$\Nat$ 的期望。

因此，显然我们不能声称自己回答了一个形而上学问题，例如\emph{哪些对象是自然数}。但这未尝不是好事。毕竟，在\olref[sfr][arith][ref]{sec}，我们曾强调，若把 $\Real$ 的定义（即 Dedekind 分割的集合）视为一种\emph{发现}而非一种方便的约定，那就错了。关键的观察在于：Dedekind 分割体现了实数的核心数学性质。此处亦然：关键的观察在于，\emph{任何} Dedekind 代数都体现了自然数的核心数学性质。（事实上，Dedekind 通过证明所有 Dedekind 代数都是\emph{同构的}（\citeyear[定理 132--3]{Dedekind1888}），有力地强调了这一点。因此，许多当代“结构主义者”将 Dedekind 引为先驱，也就不足为奇了。）

此外，我们已展示了如何将自然数的理论嵌入到一个朴素简单的集合论中；这个集合论本身仍然相当非形式化，但它（看似）并没有预设自然数是给定的。因此，我们或许正迈向实现 Dedekind 本人那雄心勃勃的计划，他曾如此解释道：

\begin{quote}
	在科学中，凡是能够证明的，都不应在未经证明时被相信。尽管这一要求看似合理，但我认为，即便是在为这门最简易科学——即，处理数论的那部分逻辑——奠基的最新方法中，它也尚未得到满足。当我说算术（代数、分析）仅仅是逻辑的一部分时，我的意思是，我认为数的概念完全独立于空间与时间的观念或直观——我更倾向于认为它是纯粹思维规律的一个直接产物。\citep[前言]{Dedekind1888}
\end{quote}

Dedekind 的大胆想法是这样的。我们刚刚展示了如何仅用（朴素）集合论来构建自然数。在\olref[sfr][arith][]{chap}，我们看到了如何在给定自然数与一些集合论的情况下构造实数。因此，或许，“算术（代数，分析）”最终仅仅是“逻辑的一部分”（在 Dedekind 对“逻辑”一词的广义理解下）。

想法就是这样。但请等一等。我们构造 Dedekind 代数（作为自然数的替代物）是以存在一个 Dedekind 无限集为条件的。（只需回顾\olref[sfr][infinite][dedekind]{thm:DedekindInfiniteAlgebra}。）除非一个 Dedekind 无限集的存在性能够通过“逻辑”或“纯粹思维规律”确立，否则这个计划就会停滞。

那么，一个 Dedekind 无限集的存在性\emph{能}被“纯粹思维规律”确立吗？以下是 Dedekind 的尝试：

\begin{quote}
  我自己的思想王国，即所有能成为我思想对象的事物的总体 $S$，是无限的。因为如果 $s$ 表示 $S$ 的一个元素，那么“$s$ 能成为我思想的对象”这一思想 $s'$，本身也是 $S$ 的一个元素。如果我们将其视为元素 $s$ 的一个映像 $\phi(s)$，则……$S$ 是 Dedekind 无限的，此即所欲证者。\citep[\S66]{Dedekind1888}
\end{quote}

在一本主要由高度严谨的数学证明构成的书中看到这些内容，是相当令人惊讶的。有两点值得说明。

第一：这个“证明”几乎不具备我们现在所认可的“数学”特征。它谈及的是心理对象（思想），并且仅仅是\emph{可能}的心理对象。

第二：至少就我们介绍 Dedekind 代数的方式而言，这个“证明”存在一个直接的技术缺陷。如果 Dedekind 的论证是成功的，它仅仅确立了存在无限多个事物（具体来说，无限多个思想）。但 Dedekind 还需要给我们一个理由，将~$S$ 视为一个单一的、拥有无限多元素的\emph{集合}，而不是将~$S$ 视为\emph{一些事物}（复数形式）。

Dedekind 未察觉到这里的空隙，这一事实可能表明，他对“总体”一词的使用并不精确对应于\emph{我们}对“集合”一词的使用。\footnote{事实上，我们也有其他理由认为确实如此；参见 \citet[p.~23]{Potter2004}。} 但这并不太令人意外。我们在前两章所推进的计划——对自然数的“构造”，并由此“构造”整数、实数和有理数——都是在朴素的意义上进行的。我们根据需要，径直取用了这个或那个集合，而没有明确规定究竟存在哪些集合、以及在何时存在的一般原则。但我们知道我们需要\emph{一些}一般原则，否则我们将陷入 Russell 悖论。

是时候超越我们的朴素性了。

\end{document}